\subsection{A valid alternative to the conservative EB gate}
\label{sec:betting-gate}

The range-dependent term in the empirical Bernstein radius is a
sufficient-bound cost, which need not be unavoidable. We therefore also
test a standard bounded-mean betting construction
\cite[Section 4]{waudby2024}. This changes the evidence mechanism while
keeping proposals and training fixed; it is not a new information score
or a new statistical test.

For one comparison with required gain $\tau\in[0,1]$, put $Z_i=D_i-\tau$,
where $D_i\in[-1,1]$ is the paired Brier difference. Our implemented
thresholds are $\tau=0$ for continuation and $\tau=0.0005$ for growth.
Fix
\[
 \Lambda_\tau=\{0.05,0.1,0.2,0.5,(1+\tau)^{-1}\},\qquad
 E_n=\frac15\sum_{\lambda\in\Lambda_\tau}
          \prod_{i=1}^n(1+\lambda Z_i).
\]
For these thresholds every $\lambda\in\Lambda_\tau$ belongs to
$[0,(1+\tau)^{-1}]$. The gate accepts when $E_n\ge1/\alpha_\mathrm{pair}$.
It computes each product in logarithms and combines them by log-sum-exp;
zero factors give a zero product. The average over the five fixed values
is essential: selecting the best product without its mixture penalty
would not have this validity statement.

\begin{proposition}[Conditional validity of the implemented betting gate]
\label{prop:betting-gate}
For the implemented thresholds $\tau\in\{0,0.0005\}$,
conditional on an arbitrary training and decision history $H$, suppose
the compared predictions and $\tau$ are fixed, the audit observations
are iid, and $\mathbb E[D_i\mid H]\le\tau$. Then
$\mathbb P(E_n\ge1/\alpha_\mathrm{pair}\mid H)
\le\alpha_\mathrm{pair}$. The same bound holds for crossing the threshold
at any audit sample size, provided the prediction pair stays fixed.
Allocation $\alpha_\mathrm{pair}=0.05/(3\cdot8)$ gives probability at
most $0.05$ of any false accepted inequality during the eight-opportunity,
three-comparison implemented controller, with fresh audit data at each
opportunity.
\end{proposition}
\begin{proof}
Write $M_{\lambda,0}=1$ and
$M_{\lambda,k}=\prod_{i=1}^k(1+\lambda Z_i)$.
Since $Z_i\ge-1-\tau$ and
$0\le\lambda\le(1+\tau)^{-1}$, every factor is nonnegative.
Conditional on $H$ and the first $k-1$ audit observations,
\[
 \mathbb E[M_{\lambda,k}\mid H,Z_1,\ldots,Z_{k-1}]
 =M_{\lambda,k-1}(1+\lambda\mathbb E[Z_k\mid H])
 \le M_{\lambda,k-1}.
\]
Thus each product and their fixed average form nonnegative
supermartingales, with initial value one. In particular
$\mathbb E[E_n\mid H]\le1$, and Markov's inequality proves the
fixed-size claim. For the threshold-crossing claim, let $T$ be the first
crossing. For every finite $m$, the stopped average $E_{T\wedge m}$ has
expectation at most one: this follows by summing its conditional
nonpositive increments up to the bounded stopping time. On $\{T\le m\}$
it is at least $1/\alpha_\mathrm{pair}$, so
$\mathbb P(T\le m\mid H)\le\alpha_\mathrm{pair}$.
Let $m\to\infty$. Finally integrate the conditional comparison bounds
over $H$ and sum them over the 24 possible tests. The three comparisons
may share the current audit sample; their independence is not needed.
\end{proof}

The proposition guarantees accepted inequalities, not the discovery of
useful candidates or high power. Products are restarted for each new
prediction pair. Reusing a product after a feedback-dependent change in
the prediction pair is outside this implementation and proof. The gate
adds $15n$ factor evaluations across the three comparisons per opportunity,
in addition to the branch fitting and prediction evaluations shared with
the EB procedure. The follow-up reports these extra operations.
